\documentclass[../../main.tex]{subfiles}
\begin{document}
\begin{oframed}
  \subfile{problem}
\end{oframed}

{\bf [解答]}

\begin{figure}[htb]
  \centering
  \begin{tikzpicture}[x=0.22cm, y=0.19cm, >=stealth, every node/.style={font=\small}]
    \coordinate (O) at (0, 0);
    \coordinate (A) at (0, 6);
    \coordinate (B) at (0, 20);
    \coordinate (P) at (11, 0);

    % 座標軸 (xz 平面)
    \draw[->, thick, gray!80!black] (-2.5, 0) -- (17.5, 0) node[right] {$x$};
    \draw[->, thick, gray!80!black] (0, -2.2) -- (0, 23.5) node[above] {$z$};
    \node[below left=1pt, font=\footnotesize] at (O) {$O$};

    % 線分 PA, PB と角 APB
    \pic[draw=red!85!black, thick, angle radius=0.85cm] {angle = B--P--A};
    \draw[line width=1.3pt, blue!85!black] (P) -- (A);
    \draw[line width=1.3pt, blue!85!black] (P) -- (B);

    % 点とラベル
    \node[circle, fill=blue!85!black, inner sep=1.3pt] at (A) {};
    \node[left=2pt, font=\footnotesize] at (A) {$6$};
    \node[above right=1pt, font=\footnotesize, blue!85!black] at (A) {$A$};

    \node[circle, fill=blue!85!black, inner sep=1.3pt] at (B) {};
    \node[left=2pt, font=\footnotesize] at (B) {$20$};
    \node[right=2pt, font=\footnotesize, blue!85!black] at (B) {$B$};

    \node[circle, fill=red!85!black, inner sep=1.4pt] at (P) {};
    \node[below=2pt, font=\footnotesize, red!85!black] at (P) {$P(x,0)$};
  \end{tikzpicture}
  \caption{$xz$ 平面における点 $A, B$ と $x$ 軸上の点 $P$}
  \label{fig:xz_plane_apb}
\end{figure}

対称性から，まず$P$が$x$軸上にある時を考える．
上図の$xz$平面において，
\begin{align}
  &\overrightarrow{PA}=
  \begin{pmatrix}-x\\6
  \end{pmatrix}&\overrightarrow{PB}=
  \begin{pmatrix}-x\\20
  \end{pmatrix}
\end{align}
だから，$0\le\angle APB<\pi/2$とあわせて，
\begin{align}
  \tan\angle APB&=\dfrac{|-6x+20x|}{\overrightarrow{PA}\cdot\overrightarrow{PB}} \\
  &=\dfrac{14x}{120+x^2} \tag{$\because x\ge0$}
\end{align}
$\tan\theta$が$0\le\theta<\pi/2$で単調増加であることから，$\angle APB\ge30^\circ$のとき，
\begin{align}
  &\tan\dfrac{\pi}{6}\le\dfrac{14x}{120+x^2} \\
  \Longleftrightarrow &x^2-14\sqrt{3}x+120 \le 0 \\
  \Longleftrightarrow &(x-10\sqrt{3})(x-4\sqrt{3})\le 0  \\
  \Longleftrightarrow&4\sqrt{3}\le x\le 10\sqrt{3}
\end{align}
したがって，$P$を$xy$平面で動かすと，
\begin{align}
  16×3\le x^2+y^2\le100×3\label{eq:1}
\end{align}
となる．従って，求める領域は
\begin{align}
  \cref{eq:1}\land(0\le x\le15)\land(0\le y\le15)
\end{align}
であり，図示して下図斜線部である．

\begin{figure}[htb]
  \centering
  \begin{tikzpicture}[scale=0.27, >=stealth, every node/.style={font=\small}]
    \pgfmathsetmacro{\rIn}{4*sqrt(3)}
    \pgfmathsetmacro{\rOut}{10*sqrt(3)}
    \pgfmathsetmacro{\yCut}{5*sqrt(3)}

    \coordinate (O) at (0, 0);
    \coordinate (Q1) at (15, \yCut);
    \coordinate (Q2) at (\yCut, 15);

    % 求める領域の斜線塗り（境界線より先に配置）
    \fill[pattern=north east lines, pattern color=blue!75!black]
    (\rIn, 0) -- (15, 0) -- (15, \yCut)
    arc[start angle=30, end angle=60, radius=\rOut]
    -- (0, 15) -- (0, \rIn)
    arc[start angle=90, end angle=0, radius=\rIn] -- cycle;

    % 座標軸
    \draw[->, thick, gray!80!black] (-2.0, 0) -- (20.8, 0) node[right] {$x$};
    \draw[->, thick, gray!80!black] (0, -1.8) -- (0, 20.5) node[above] {$y$};
    \node[below left=1pt, font=\footnotesize] at (O) {$O$};

    % 外円（領域外の部分）と正方形の補助線・垂線
    \draw[dashed, thick, blue!70!black] (\rOut, 0) arc[start angle=0, end angle=30, radius=\rOut];
    \draw[dashed, thick, blue!70!black] (Q2) arc[start angle=60, end angle=90, radius=\rOut];
    \draw[densely dotted, thick, gray!70!black] (Q1) -- (15, 15) -- (Q2);
    \draw[densely dotted, thick, gray!70!black] (Q1) -- (0, \yCut);
    \draw[densely dotted, thick, gray!70!black] (Q2) -- (\yCut, 0);

    % 面積計算用の扇形分割線と偏角
    \draw[dashed, thick, red!80!black] (O) -- (Q1);
    \draw[dashed, thick, red!80!black] (O) -- (Q2);
    \draw[thick, red!80!black] (3.6, 0) arc[start angle=0, end angle=30, radius=3.6];
    \node[text=red!80!black, fill=white, inner sep=0.8pt, font=\scriptsize] at (15:4.9) {$\frac{\pi}{6}$};
    \draw[thick, red!80!black] (30:4.2) arc[start angle=30, end angle=60, radius=4.2];
    \node[text=red!80!black, fill=white, inner sep=0.8pt, font=\scriptsize] at (45:5.5) {$\frac{\pi}{6}$};

    % 求める領域の輪郭
    \draw[line width=1.4pt, blue!85!black]
    (\rIn, 0) -- (15, 0) -- (15, \yCut)
    arc[start angle=30, end angle=60, radius=\rOut]
    -- (0, 15) -- (0, \rIn)
    arc[start angle=90, end angle=0, radius=\rIn] -- cycle;

    % 交点・端点と軸上のラベル
    \node[circle, fill=red!85!black, inner sep=1.3pt] at (Q1) {};
    \node[circle, fill=red!85!black, inner sep=1.3pt] at (Q2) {};
    \node[above right=2pt, font=\footnotesize, text=red!85!black, fill=white, inner sep=0.8pt] at (Q1) {$(15,5\sqrt{3})$};
    \node[above right=3.5pt, font=\footnotesize, text=red!85!black, fill=white, inner sep=0.8pt] at (Q2) {$(5\sqrt{3},15)$};

    \node[below=2pt, font=\footnotesize] at (\rIn, 0) {$4\sqrt{3}$};
    \node[below left=1pt, font=\footnotesize] at (15, 0) {$15$};
    \node[below right=1pt, font=\footnotesize] at (\rOut, 0) {$10\sqrt{3}$};

    \node[below left=0pt, font=\footnotesize] at (0, \rIn) {$4\sqrt{3}$};
    \node[above left=0pt, font=\footnotesize] at (0, \yCut) {$5\sqrt{3}$};
    \node[left=2pt, font=\footnotesize] at (0, 15) {$15$};
    \node[left=2pt, font=\footnotesize] at (0, \rOut) {$10\sqrt{3}$};
  \end{tikzpicture}
  \caption{$xy$ 平面上で条件を満たす点 $P$ の領域}
  \label{fig:xy_region_apb}
\end{figure}

この面積$S$は
\begin{align}
  S&=\dfrac{1}{2}\dfrac{\pi}{6}(10\sqrt{3})^2+15×5\sqrt{3}-\dfrac{1}{2}\dfrac{\pi}{2}(4\sqrt{3})^2 \\
  &=13\pi+75\sqrt{3}
\end{align}
である．$\cdots$(答)

\newpage
\end{document}
